\section{Borcherds-product expansion: \texorpdfstring{$\phi_{0,1}$}{phi01}, the
\texorpdfstring{$\Delta_5$}{Delta5} infinite product, and the
Gritsenko--Cl\'ery eight-row catalogue}
\label{app:borcherds-product-expansion}

The Fourier expansion of the
weight-zero index-one weak Jacobi form $\phi_{0,1}$, the
Borcherds--Gritsenko--Nikulin product expansion of the Igusa cusp form
$\Delta_5$ in the type-II cusp chamber, the theta-constant normalization
constant $64=2^{6}$, the doubling identity that turns the monic Borcherds
product into the denominator of the Lorentzian Borcherds--Kac--Moody
algebra $\autcor$, and the full eight-row Gritsenko--Cl\'ery catalogue of
diagonal-divisor Siegel modular forms are the automorphic data at the
type-II cusp.  Every displayed numerical Fourier
coefficient is produced by exact rational expansion of
the theta quotient; every identity is
attributed to author, year, and theorem of the primary literature.
Orientation data, Pfaffian wall normal forms, and compact Hall input do
not enter these automorphic identities; they live entirely on the
automorphic / target side, prior to the geometric content of (D0),
(O1), (O1)\(^+\), and (O2).

Three names for $\Delta_5$ appear and are kept distinct: the
theta-constant automorphic section $\mathcal D_X=\Delta_5$ (View \textcircled{\scriptsize 1}); the
Borcherds--Kac--Moody denominator
$\operatorname{den}(\autcor)=64^{-1}\Delta_5(2Z)$
(View \textcircled{\scriptsize 2}); and the protected Pfaffian
$\operatorname{Pf}_{\mathrm{prot}}(\mathfrak D_X^{\mathrm{DI}})$ of the
pro-object $\mathfrak D_X^{\mathrm{DI}}$ (View \textcircled{\scriptsize 3},
after the retained D0-HN datum, retained quotient-orientation datum,
chosen Weyl lifts, retained \((O2_{\mathrm{atlas}})\) wall datum, and
finite geometric Pfaffian datum \((P_{\mathrm{fin}})\) of
Appendix~\ref{app:obstruction-discharges}).  The first and third
sections coincide as weight-five automorphic sections; the second
coincides with their Borcherds denominator after the normalization
\(D_5=64^{-1}\Delta_5\) and the pullback \(Z\mapsto2Z\).  The
orientation character that distinguishes
$\Delta_5$ from $-\Delta_5$ enters only at View \textcircled{\scriptsize 3}.

\subsection*{B.1 The weak Jacobi form \texorpdfstring{$\phi_{0,1}$}{phi01}}

The K3 elliptic genus is the modular input.  Eichler--Zagier
\cite[\S2.2 and Theorem~9.3]{EZ:1} fix the weight-zero index-one
weak-Jacobi normalization
\[
\phi_{0,1}(\tau,z)\in J_{0,1}^{\mathrm w},
\qquad
\phi_{0,1}(\tau,z)=\sum_{n\ge0,\;l\in\Z}f(n,l)\,q^n r^l,
\qquad
q=e^{2\pi i\tau},\ r=e^{2\pi iz}.
\]
The relation to the K3 right-moving elliptic genus is
$Z_{\mathrm{K3}}(\tau,z)=2\phi_{0,1}(\tau,z)$ \cite[\S2.2]{EZ:1};
the half is the Borcherds-lift representative for the Borcherds product.
The cusp normalization is
\[
\phi_{0,1}(\tau,z)=r^{-1}+10+r+O(q).
\]
The Jacobi weak-form support is
$f(n,l)=0$ whenever $4n-l^{2}<-1$, with isolated extreme entries
$f(0,\pm1)=1$ on the discriminant hyperbola $4n-l^{2}=-1$.
\emph{(proved; \cite[\S2.2, Theorem~9.3]{EZ:1};
\cite[Example~2.4]{GNII}.)}

\paragraph{Verified Fourier table.}
The exact coefficients $f(n,l)$ for $n\le2$ are the rational integers
computed by the theta-quotient identity
\(\phi_{0,1}=4\sum_{i=2,3,4}\theta_i(\tau,z)^2/\theta_i(\tau,0)^2\) of
\cite[\S2.2]{EZ:1} by rational arithmetic:
\[
\begin{array}{c|ccccccc}
& l=-3 & l=-2 & l=-1 & l=0 & l=1 & l=2 & l=3 \\ \hline
n=0 & 0 & 0 & 1 & 10 & 1 & 0 & 0 \\
n=1 & 0 & 10 & -64 & 108 & -64 & 10 & 0 \\
n=2 & 1 & 108 & -513 & 808 & -513 & 108 & 1
\end{array}
\]
\emph{(computed; cross-checked with \cite[Example~2.4]{GNII} and
\cite[Theorem~4.1]{GN}.)}
The finite table
\(\texttt{certificates/jacobi/k3e\_first\_relation\_window\_coefficients}\)
records this leading normalization together with the checked evenness
and discriminant-dependence rows used in the first relation-closed
window.
The negative entries reflect the index-one Jacobi-form theta-decomposition
$\phi_{0,1}=\theta_{1,0}h_0+\theta_{1,1}h_1$ with theta-coefficient
expansions $h_0=2(45 E_4 E_{4,1}-50 E_6 E_{6,1})/\eta^{24}+\cdots$ in
\cite[\S2.2]{EZ:1}; the integer values are not target signs but signed
multiplicities of the Borcherds source.

\paragraph{Fourier identities.}
\[
\begin{aligned}
f(0,0)&=10,&\qquad f(0,\pm1)&=1,&\qquad f(0,l)&=0\ \ (|l|\ge2),\\
f(1,0)&=108,&\qquad f(1,\pm1)&=-64,&\qquad f(1,\pm2)&=10,\\
f(2,0)&=808,&\qquad f(2,\pm1)&=-513,&\qquad f(2,\pm2)&=108,\\
&&\qquad f(2,\pm3)&=1,&\qquad f(1,\pm3)&=0.
\end{aligned}
\]
The diagonal entry $f(0,0)=10$ controls the Borcherds-lift weight
\(\operatorname{wt}(\Delta_5)=f(0,0)/2=5\); the entries $f(0,\pm1)=1$
control the simple type-II divisor orders
$d_{\delta_i}^{\mathrm{aut}}=1$; the entry $f(1,1)=-64$ controls the
multiplicity of the first timelike imaginary root, computed and used
below.

\subsection*{B.2 The cusp chamber and the effective semigroup}

Set $Z=\bigl(\begin{smallmatrix}z_1&z_2\\z_2&z_3\end{smallmatrix}\bigr)
\in\Hpl_2$, $q=e^{2\pi iz_1}$, $r=e^{2\pi iz_2}$, $s=e^{2\pi iz_3}$.
The standard type-II cusp polarization is the ordered limit
\[
0<|s|\ll|q|\ll|r|^{-1}\ll1,
\]
which determines the effective semigroup
\[
\Gamma_{\mathrm{eff}}
=\{(n,l,m)\in\Z^3\mid m>0,\ n\ge0\}
\cup\{(n,l,0)\mid n>0\}
\cup\{(0,l,0)\mid l<0\}.
\]
Inside the Gram-index lattice $\Gamma_{\mathrm{ind}}=\Z^3$, the
\emph{active support}
\[
\Gamma_{\mathrm{act}}
:=\{(n,l,m)\in\Gamma_{\mathrm{eff}}\mid f(nm,l)\ne0\}
\]
is the set of triples that contribute a visible product factor; the
remaining triples in $\Gamma_{\mathrm{eff}}$ have zero exponent and fix
ordering only.  The signed Gram pairing
\((\alpha(n,l,m),\alpha(n,l,m))=-2(4nm-l^{2})\)
under
\[
\alpha:(n,l,m)\longmapsto 2nf_2-lf_3+2mf_{-2}\in\La^{2,1}_{II}
\]
restricts to the Lorentzian root lattice of $\autcor$.
The finite packet
\(\texttt{certificates/lattice/product\_chamber\_root\_map}\)
gives a computed check of the three sector definitions of
\(\Gamma_{\mathrm{eff}}\), their pairwise closure under addition, the
active rows selected by \(f(nm,l)\ne0\), the inactive zero-exponent rows,
and the identity
\[
\alpha(n,l,m)=n\delta_1+m\delta_2+(n+m-l)\delta_3
\]
used in the type-II chamber.  This is target-side lattice arithmetic; it
does not construct a compact \(K3\times E\) source, a Pfaffian
orientation, an \(O2\) wall atlas, a mirror discriminant, or a protected
trace.

\begin{proposition}[Images of the three simple type-II walls]
\label{prop:type-ii-wall-images}
\textnormal{[computed]}
Let
\[
\gamma_{\mathrm{wall}}:
\{\delta_1,\delta_2,\delta_3\}\longrightarrow \Gamma_{\mathrm{eff}}
\subset\Z^3
\]
be the inverse of the product-chamber root map
\[
\alpha(n,l,m)
=n\delta_1+m\delta_2+(n+m-l)\delta_3
\]
on the three simple type-II roots.  Then
\[
\gamma_{\mathrm{wall}}(\delta_1)=(1,1,0),\qquad
\gamma_{\mathrm{wall}}(\delta_2)=(0,1,1),\qquad
\gamma_{\mathrm{wall}}(\delta_3)=(0,-1,0).
\]
Equivalently, under
\[
\alpha(n,l,m)=2nf_2-lf_3+2mf_{-2},
\]
these triples map respectively to
\[
2f_2-f_3,\qquad 2f_{-2}-f_3,\qquad f_3 .
\]
This is a target-side wall-image computation.  It does not construct
compact \(K3\times E\) representatives, retained type-II wall objects,
wall charts, Pfaffian signs, protected traces, or a Hall carrier.
\end{proposition}

\begin{proof}
For a triple \((n,l,m)\), the coefficient of
\((\delta_1,\delta_2,\delta_3)\) in the type-II chamber is
\[
(n,\ m,\ n+m-l).
\]
For \(\delta_1\) this coefficient vector must be \((1,0,0)\), hence
\[
n=1,\qquad m=0,\qquad 1-l=0,
\]
so \(l=1\) and \(\gamma_{\mathrm{wall}}(\delta_1)=(1,1,0)\).
For \(\delta_2\) the coefficient vector is \((0,1,0)\), hence
\[
n=0,\qquad m=1,\qquad 1-l=0,
\]
so \(l=1\) and \(\gamma_{\mathrm{wall}}(\delta_2)=(0,1,1)\).
For \(\delta_3\) the coefficient vector is \((0,0,1)\), hence
\[
n=0,\qquad m=0,\qquad -l=1,
\]
so \(l=-1\) and \(\gamma_{\mathrm{wall}}(\delta_3)=(0,-1,0)\).
Substituting the three triples into
\(\alpha(n,l,m)=2nf_2-lf_3+2mf_{-2}\) gives
\[
2f_2-f_3,\qquad -f_3+2f_{-2},\qquad f_3,
\]
which are precisely \(\delta_1,\delta_2,\delta_3\).
\end{proof}

The finite packet
\[
\texttt{certificates/lattice/type\_ii\_wall\_images}
\]
verifies
\(\texttt{TYPE\_II\_WALL\_IMAGES\_VERIFIED}\): the three inverse
images above have zero root-map defect and lie in the product chamber
sectors.  The packet records compact source representatives, retained
wall objects, wall charts, Pfaffian signs, protected traces, and Hall
carrier data as excluded scalar substitutes.

\subsection*{B.3 The Borcherds product formula for
\texorpdfstring{$\Delta_5$}{Delta5}}

\begin{theorem}[Borcherds--Gritsenko--Nikulin product expansion]
\label{thm:borcherds-product-delta5}
On the type-II cusp polarization $0<|s|\ll|q|\ll|r|^{-1}\ll1$, the
weight-five Igusa cusp form admits the absolutely convergent infinite
product
\[
\Delta_5(Z)
=64\,q^{1/2}r^{1/2}s^{1/2}
\!\!\prod_{(n,l,m)\in\Gamma_{\mathrm{eff}}}\!\!
(1-q^n r^l s^m)^{f(nm,l)},
\]
with theta-leading coefficient
\[
[q^{1/2}r^{1/2}s^{1/2}]\Delta_5\;=\;64\;=\;2^{6}.
\]
The monic product is $D_5:=64^{-1}\Delta_5$; the constant $64$ is the
theta-constant Igusa normalization, not a separately chosen scalar.
\emph{(proved; \cite[Theorem~10.1]{Borcherds95},
\cite[Theorem~2.1, Example~2.4, equations (2.15)--(2.16)]{GNII},
\cite[Theorem~4.1]{GN}.)}
\end{theorem}

\begin{proof}
The general Borcherds singular theta-lift sends a weakly holomorphic
vector-valued modular form of negative weight for the Weil
representation of a lattice of signature \((2,n)\) to a meromorphic
orthogonal automorphic form realized as an infinite product
\cite[Theorem~10.1]{Borcherds95}; the regularisation and convergence on
divisors with singularities in higher rank is in
\cite[\S\S6--13]{B}.  The genus-two specialization at the
Siegel cusp uses the exterior-square isogeny
$\PSp_4(\R)\simeq\SO(3,2)_+$ and the
Eichler--Zagier dictionary $J^{\mathrm w}_{0,1}\leftrightarrow$
vector-valued Mp$_{2}$-modular forms \cite[\S5]{EZ:1};
the ordered limit $0<|s|\ll|q|\ll|r|^{-1}\ll1$ selects the type-II cusp
\cite[\S2, Example~2.4]{GNII}.  Gritsenko--Nikulin \cite[(2.15)--(2.16)]{GNII}
write the resulting product in the unnormalized form
\[
D_5(Z)
=q^{1/2}r^{1/2}s^{1/2}
\prod_{(n,l,m)\in\Gamma_{\mathrm{eff}}}(1-q^n r^l s^m)^{f(nm,l)},
\]
and identify $D_5$ with $64^{-1}$ times the theta-constant cusp form
$\Delta_5$ of Igusa \cite[Theorem~4.1]{GN}, \cite{Igusa1962},
\cite{Igusa1964II}; the leading constant $2^{6}=64$ is the explicit
theta-product coefficient.
\end{proof}

\paragraph{The constant $2^{6}$ as a theta-product leading coefficient.}
Igusa's classical product formula displays $\Delta_5$ as the product of
the ten even theta constants of genus two,
\[
\Delta_5(Z)
=\prod_{\substack{a,b\in(\Z/2\Z)^2\\{}^{t}\!ab\equiv0\pmod2}}
\nu_{a,b}(Z),
\qquad
\nu_{a,b}(Z)
=\sum_{l\in\Z^{2}}\!\!\exp\!\bigl(\pi i(Z[l+\tfrac12 a]+{}^{t}\!bl)\bigr),
\]
\cite{Igusa1962}, \cite[Theorem~4.1]{GN}, \cite[\S6]{FREITAG:1}.
In the type-II cusp polarization, with $z_2$ the off-diagonal modulus,
the four characteristics $a=(0,0)$ start with $1$, the two
characteristics $a=(1,0)$ start with $2q^{1/8}$, the two
characteristics $a=(0,1)$ start with $2s^{1/8}$, and the two
characteristics $a=(1,1)$ start with $2q^{1/8}r^{1/4}s^{1/8}$; the
total leading monomial is $q^{1/2}r^{1/2}s^{1/2}$ with constant
$1\cdot 2^{2}\cdot 2^{2}\cdot 2^{2}=2^{6}=64$
\emph{(proved; \cite{Igusa1962}, \cite[Theorem~4.1]{GN}.)}
This is not a normalization choice: the constant $64$ is read off the
theta product, with no scalar freedom.  Replacing $\Delta_5$ by
$c\Delta_5$, $c\in\C^{\times}$, would change neither the Maass character
\(\nu_{\Delta_5}\) nor the divisor orders
\(d_{\delta}^{\mathrm{aut}}\) of \cite[Theorem~4.1]{GN}.

\subsection*{B.4 Sample exponent verifications}

The exponent $f(nm,l)$ in the product formula is read directly from the
Jacobi-coefficient table of B.1.  The following sample is verified
exactly by the displayed coefficient table:

\begin{center}
\small
\begin{tabular}{c|cccccc|c}
\hline
factor $(n,l,m)$ & \multicolumn{6}{c|}{Lorentzian Gram label
$\alpha=2nf_{2}-lf_{3}+2mf_{-2}$} & exponent $f(nm,l)$ \\
\hline
$(1,1,0)$ & \multicolumn{6}{c|}{$\delta_{1}=2f_{2}-f_{3}$} & $f(0,1)=1$ \\
$(0,1,1)$ & \multicolumn{6}{c|}{$\delta_{2}=2f_{-2}-f_{3}$} & $f(0,1)=1$ \\
$(0,-1,0)$ & \multicolumn{6}{c|}{$\delta_{3}=f_{3}$} & $f(0,-1)=1$ \\
$(0,0,m)\ (m\ge1)$ & \multicolumn{6}{c|}{cusp ray
$2mf_{-2}$} & $f(0,0)=10$ \\
$(1,0,1)$ & \multicolumn{6}{c|}{$2f_{2}+2f_{-2}$ (timelike,
$(\alpha,\alpha)=-8$)} & $f(1,0)=108$ \\
$(1,1,1)$ & \multicolumn{6}{c|}{$2f_{2}-f_{3}+2f_{-2}$ (timelike,
$(\alpha,\alpha)=-6$)} & $f(1,1)=-64$ \\
$(1,2,1)$ & \multicolumn{6}{c|}{$2f_{2}-2f_{3}+2f_{-2}$ (isotropic,
$(\alpha,\alpha)=0$)} & $f(1,2)=10$ \\
$(1,3,1)$ & \multicolumn{6}{c|}{$2f_{2}-3f_{3}+2f_{-2}$
(spacelike, $(\alpha,\alpha)=10$)} & $f(1,3)=0$ \\
$(2,1,1)$ & \multicolumn{6}{c|}{$4f_{2}-f_{3}+2f_{-2}$
($(\alpha,\alpha)=-14$)} & $f(2,1)=-513$ \\
$(1,2,2)$ & \multicolumn{6}{c|}{$2f_{2}-2f_{3}+4f_{-2}$
($(\alpha,\alpha)=-12$)} & $f(2,2)=108$ \\
\hline
\end{tabular}
\end{center}

The triple $(1,1,1)$ gives the first timelike imaginary root: its
exponent $f(1,1)=-64$ is the signed superdimension of the root space
$\mathfrak g_{\Delta_5,\alpha(1,1,1)}$ on the active support, and is the
content of the lattice cross-check
$f(1,1)=29-93=-64$ from the
free-Lie/imaginary-simple split computed by
the Gritsenko--Nikulin target presentation.  The triple $(1,3,1)$ has
exponent zero because $f(1,3)=0$: this triple lies in
$\Gamma_{\mathrm{eff}}\setminus\Gamma_{\mathrm{act}}$ and contributes
ordering, not a visible factor.  These are signed superdimensions, not
positive dimensions; negative entries record an odd excess in the
Borcherds source, and only after the (O1), (O1)\(^+\), (O2) orientation
data become available do they translate to operator parities.

\paragraph{Theta-leading coefficient.}
The coefficient \(64\) is the coefficient of the leading monomial of
\(\Delta_5\):
\[
\Delta_5(Z)=64\,q^{1/2}r^{1/2}s^{1/2}\bigl(1+\cdots\bigr).
\]
After the leading monomial is divided out, the first \(qrs\)-coefficient
of \(D_5=64^{-1}\Delta_5\) is \(93\), not \(64\).  The factor
\((1-qrs)^{-64}\) contributes \(64\,qrs\), and the remaining
root factors contribute the other \(29\,qrs\); this is the identity
\[
[qrs]\,D_5/(qrs)^{1/2}=93,\qquad 29-93=-64=f(1,1).
\]
\emph{(computed by exact rational expansion of the Borcherds product.)}

\subsection*{B.5 The doubling identity
\texorpdfstring{$\operatorname{den}(\autcor)=64^{-1}\Delta_5(2Z)$}
{den(g\_Delta5)=64\^{}-1 Delta5(2Z)}}

\begin{lemma}[Two exponential normalizations]
\label{lem:appB-doubling}
The product chamber convention writes monomials as $q^{n}r^{l}s^{m}$
with $q^{n}r^{l}s^{m}=\exp(-\pi i(\alpha(n,l,m),z))$, while the standard
denominator-side convention of a Borcherds--Kac--Moody algebra writes
formal characters as $\exp(-2\pi i(\alpha,z))$.  Replacing $Z$ by $2Z$
converts the first character to the second, so
\[
\operatorname{den}(\autcor)
\;=\;D_{5}(2Z)
\;=\;64^{-1}\Delta_{5}(2Z).
\]
\emph{(proved; \cite[Proposition~3.1]{GN}, \cite[Section~5.1.1]{GNII}
together with Theorem~\ref{thm:borcherds-product-delta5}.)}
\end{lemma}

\begin{proof}
The product side of Theorem~\ref{thm:borcherds-product-delta5} gives
\(D_{5}(Z)=q^{1/2}r^{1/2}s^{1/2}\prod(1-q^{n}r^{l}s^{m})^{f(nm,l)}\)
with leading monomial $q^{1/2}r^{1/2}s^{1/2}$.  Doubling
$Z\mapsto 2Z$ replaces $q^{n}r^{l}s^{m}$ by $q^{2n}r^{2l}s^{2m}$ and
hence $\exp(-\pi i(\alpha,z))$ by $\exp(-2\pi i(\alpha,z))$.  The
Lorentzian Borcherds--Kac--Moody algebra $\autcor$ of
Gritsenko--Nikulin \cite[Sections~3--4, Proposition~3.1]{GN} carries the
denominator identity
\[
\operatorname{den}(\autcor)
=\sum_{w\in W^{(2)}(\La^{2,1}_{II})}\det(w)\,
\Bigl(e^{-2\pi i(w\rho,z)}
-\sum_{a}m(a)\,e^{-2\pi i(w(\rho+a),z)}\Bigr)
\]
on the imaginary-positive cone, where \(m(a)\) is determined by
the negative-norm additive coefficients and by
\(1-\sum_{t\ge1}m(ta)q^t=\prod_{t\ge1}(1-q^t)^{\tau(ta)}\) on isotropic
rays.  The right-hand side is equal to
$D_{5}(2Z)$ in the doubled exponential.  Equivalently
$\operatorname{den}(\autcor)=64^{-1}\Delta_{5}(2Z)$.
\end{proof}

The factor $64$ in $\operatorname{den}(\autcor)=64^{-1}\Delta_{5}(2Z)$
is the same theta-constant $2^{6}$ from the leading-coefficient
calculation, transported through
$Z\mapsto 2Z$; it is not a separate normalization, an integration
constant, or a free scalar.  The three identifications
\(\mathcal D_{X}=\Delta_{5}\),
\(\operatorname{den}(\autcor)=64^{-1}\Delta_{5}(2Z)\),
\(\operatorname{Pf}_{\mathrm{prot}}(\mathfrak D_{X}^{DI})\) are kept
distinct; the doubling lemma supplies the denominator-character
normalization.

\subsection*{B.6 The eight-row Gritsenko--Cl\'ery catalogue}

Gritsenko and Cl\'ery \cite[Theorem~1.4]{GC} classify all Siegel modular
forms of genus two, with character or multiplier system on a
Hecke-type congruence subgroup $\Gamma_{t}(N)\subset\operatorname{Sp}(4,\Q)$,
whose divisor is supported on the $\Gamma_{t}(N)$-translates of the
diagonal Humbert surface
$\mathcal H_{\mathrm{diag}}=\{Z=\operatorname{diag}(a,b):a,b\in\mathbb H_{1}\}$
with multiplicity one.  The classification produces exactly eight rows.
\emph{(proved; \cite[Theorem~1.4 and Theorem~3.1]{GC}.)}

The rows are:

\begin{center}
\renewcommand{\arraystretch}{1.16}
\small
\begin{tabular}{c|c|c|c|c|c|c}
\hline
$j$ & $F_{j}$ & $(t,N)$ & $\Gamma_{j}$ &
$\operatorname{wt}(F_{j})$ & multiplier $\nu_{j}$ & $c_{j}(0)$ \\ \hline
$1$ & $\Delta_{5}$ & $(1,1)$ & $\Gamma_{1}=\operatorname{Sp}(4,\Z)$ & $5$ &
$v_{\eta}^{12}\times v_{H}$ & $10$ \\
$2$ & $\Delta_{2}$ & $(2,1)$ & $\Gamma_{2}^{+}$ & $2$ &
$v_{\eta}^{6}\times v_{H}$ & $4$ \\
$3$ & $\Delta_{1}$ & $(3,1)$ & $\Gamma_{3}^{+}$ & $1$ &
$v_{\eta}^{4}\times v_{H}$ & $2$ \\
$4$ & $\Delta_{1/2}$ & $(4,1)$ & $\Gamma_{4}^{+}$ & $\tfrac12$ &
$v_{\eta}^{3}\times v_{H}$ (degree 8) & $1$ \\
$5$ & $\nabla_{3}$ & $(1,2)$ & $\Gamma_{0}^{(2)}(2)$ & $3$ &
$\chi_{2}^{(2)}\times v_{H}$ & $6$ \\
$6$ & $\nabla_{2}$ & $(1,3)$ & $\Gamma_{0}^{(2)}(3)$ & $2$ &
$\chi_{2}^{(3)}\times v_{H}$ & $4$ \\
$7$ & $\nabla_{3/2}$ & $(1,4)$ & $\Gamma_{0}^{(2)}(4)$ & $\tfrac32$ &
$\chi_{4}^{(4)}\times v_{H}$ (order 4) & $3$ \\
$8$ & $Q_{1}$ & $(2,2)$ & $\Gamma_{2}(2)$ & $1$ &
$\chi_{4}^{(2)}\times v_{H}$ & $2$ \\
\hline
\end{tabular}
\end{center}

In the displayed \(F_j\)-order the weights and constant Fourier
coefficients are
\[
(5,2,1,\tfrac12,3,2,\tfrac32,1),\qquad
(10,4,2,1,6,4,3,2).
\]
In the input-pair order
\[
(1,1),(2,1),(1,2),(3,1),(1,3),(4,1),(1,4),(2,2)
\]
the same data read
\[
(5,2,3,1,2,\tfrac12,\tfrac32,1),\qquad
(10,4,6,2,4,1,3,2).
\]
Each row satisfies the Borcherds--Gritsenko universal weight identity
\[
\operatorname{wt}(F_{j})\;=\;\frac{c_{j}(0)}{2}
\;=\;\kappa_{\mathrm{BKM}}(F_{j}),
\qquad j=1,\ldots,8 .
\]
\emph{(proved; \cite{Borcherds95},
\cite[Theorem~3.1]{GC}; the $\kappa_{\mathrm{BKM}}$ relabelling is the
universal Borcherds-weight identity used in the Calabi--Yau
quantum-groups $\kappa$-stratification.)}

The half-integer weights $\tfrac12$ (row $4$) and $\tfrac32$ (row $7$)
are realized by multiplier systems $v_{\eta}^{k}\times v_{H}$ on the
ambient paramodular cover, not by passing to a metaplectic double
cover.  The multiplier $v_{\eta}$ is the Dedekind eta multiplier on
$\operatorname{SL}_{2}(\Z)$; the multiplier $v_{H}$ is the Heisenberg
multiplier on the Heisenberg lift of $\operatorname{Sp}(4,\Z)$; and
$v_{\eta}^{3}\times v_{H}$ has order $8$ on the row-$4$ cover, while
$\chi_{4}^{(4)}\times v_{H}$ has order $4$ on the row-$7$ cover.
\emph{(proved; \cite[Theorem~3.1]{GC}.)}

\paragraph{Denominator content of the rows.}
Rows $1$--$4$ carry Gritsenko--Nikulin denominator theorems
\cite[Theorem~2.1, equations (2.11), (2.20), (2.21)]{GNII}; in the
$e_{t}^{(n,l,m)}=q^{n}r^{l}s^{tm}$ exponential normalization of
\cite[\S5.1.1]{GNII}, the products are
\[
\begin{aligned}
\Delta_{2}(Z)&=q^{1/4}r^{1/2}s^{1/2}\!\!
\prod_{(n,l,m)\in\Gamma_{2}^{\mathrm{prod}}}\!\!
(1-q^{n}r^{l}s^{2m})^{f_{2}(nm,l)},\\
\Delta_{1}(Z)&=q^{1/6}r^{1/2}s^{1/2}\!\!
\prod_{(n,l,m)\in\Gamma_{3}^{\mathrm{prod}}}\!\!
(1-q^{n}r^{l}s^{3m})^{f_{3}(nm,l)},\\
\Delta_{1/2}(Z)&=q^{1/8}r^{1/2}s^{1/2}\!\!
\prod_{(n,l,m)\in\Gamma_{4}^{\mathrm{prod}}}\!\!
(1-q^{n}r^{l}s^{4m})^{f_{4}(nm,l)},
\end{aligned}
\]
with cusp Jacobi inputs $\phi_{0,2},\phi_{0,3},\phi_{0,4}$ supplied
explicitly by \cite[\S5]{GNII}.  Rows $5$--$8$ carry Cl\'ery--Gritsenko
products on the corresponding congruence covers, with cusp Jacobi inputs
$\phi_{j}$ as in \cite[Theorem~3.1]{GC}, and the additional
denominator data (Weyl chamber, simple real roots, reflection character,
Weyl alternating sum) supplied by Govindarajan's twisted-CHL theorem
\cite{Gov11BKM}, \cite{GK09CHL}, \cite{CD09BKM} for $F_{5}=\nabla_{3}$,
$F_{6}=\nabla_{2}$, $F_{7}=\nabla_{3/2}$, $F_{8}=Q_{1}$.  For $F_{7}$ the
denominator theorem lives on the order-four multiplier cover
$\Gamma_{0}^{(2)}(4)\langle\chi_{4}^{(4)}\times v_{H}\rangle$.
\emph{(proved; \cite{Gov11BKM},
\cite[Theorem~2.1]{GNII}, \cite[Theorem~3.1]{GC}.)}

\paragraph{Level separation.}
The eight-row catalogue is a list of \emph{automorphic products plus
denominator theorems on the rows where the data have been supplied}.
It is not a list of compact Hall sources, reduced curve-counting
backgrounds, or chiral hosts.  For a row \(j\), the automorphic product
\(F_j\) acquires a compact physical host only after separate choices of
an automorphic line, a denominator algebra, a scalar trace, a chiral
source, and a Hall realization; these choices are row-dependent.  The
row $j=1$
($F_{1}=\Delta_{5}$) is the only row with a proved reduced CY/DT scalar
host (Oberdieck--Pixton, Oberdieck--Pandharipande:
$Z^{X}_{\mathrm{OP}}=-4096\Delta_{5}^{-2}$); even row $1$ has open
compact Hall/chiral host at the primitive-recognition level, governed
by \((S1)\)--\((S10)\).

\subsection*{B.7 The Borcherds-weight coordinate \texorpdfstring{$\kappa_{\mathrm{BKM}}=5$}{kappa\_BKM=5}}

The universal Borcherds-weight identity
\(\kappa_{\mathrm{BKM}}(\Phi_{N})=c_{N}(0)/2\) is the row-by-row
specialization of the Borcherds singular-theta-lift weight formula
\cite[Theorem~10.1]{Borcherds95} together with the Gritsenko--Nikulin
unification \cite[\S2]{GN}, \cite{Gritsenko99}.  The relevant row is
$j=1$, with
\[
\kappa_{\mathrm{BKM}}(\Delta_{5})=\operatorname{wt}(\Delta_{5})=5
=\frac{c_{1}(0)}{2}=\frac{f(0,0)}{2}=\frac{10}{2}.
\]
\emph{(proved; \cite[Theorem~10.1]{Borcherds95},
\cite[\S2.2]{EZ:1}, \cite[Theorem~3.1]{GC}.)}

This is the Borcherds-weight invariant: the same integer
$\kappa_{\mathrm{BKM}}=5$ identifies the \(\Delta_5\) row inside the
companion Calabi--Yau quantum-groups $\kappa$-stratification, and it
is the entry that distinguishes
$\kappa_{\mathrm{BKM}}=\operatorname{wt}(\Delta_{5})$ (the
Igusa cusp-form weight) from the unrelated quantities
$\kappa_{\mathrm{ch}}^{\mathrm{Hodge}}=0$,
$\kappa_{\mathrm{ch}}^{\mathrm{Heis}}=3$, $\kappa_{\mathrm{cat}}=0$,
$\kappa_{\mathrm{fiber}}=24=\chi_{\mathrm{top}}(\mathrm{K3})$.  The
additive formula $\kappa_{\mathrm{BKM}}=\kappa_{\mathrm{ch}}+
\chi(\mathcal O_{\mathrm{fiber}})$ would identify
$24=\chi_{\mathrm{top}}(\mathrm{K3})$ with
$2=\chi(\mathcal O_{\mathrm{K3}})$.  The Heisenberg reading gives an
accidental equality at \(N=1\), but it does not extend to the CHL
weights \(c_N(0)/2=(5,3,2,\tfrac32,1)\); the appendix
asserts $\kappa_{\mathrm{BKM}}=5$ as the row entry of the
Borcherds-weight identity, not as the value of any additive formula.

\subsection*{B.8 Normalizations and residual obstructions}

The following normalizations fix the Borcherds product without residual
scalar, chamber, multiplier, or coefficient ambiguity.

\begin{itemize}
\item \emph{Theta-constant normalization.}  The constant \(64\) is not
a free scalar on \(\Delta_{5}\).  The identity \(64=2^{6}\) is the
explicit theta-product leading coefficient
\cite{Igusa1962}, \cite[Theorem~4.1]{GN}; with the theta normalization
fixed, no residual scalar freedom remains.  Cross-checked against
the consistency
$f(1,1)=-64$ and the \(qrs\)-coefficient identity \(c_{\Delta}(1,1,1)=64\)
\cite[Theorem~4.1]{GN}.
\item \emph{Cusp chamber definition \(\Gamma_{\mathrm{eff}}\).}
The semigroup $\Gamma_{\mathrm{eff}}$ is the ordered-limit semigroup of
the type-II cusp polarization $0<|s|\ll|q|\ll|r|^{-1}\ll1$
\cite[Example~2.4]{GNII}; not a free choice.  Triples in
$\Gamma_{\mathrm{eff}}\setminus\Gamma_{\mathrm{act}}$ have zero exponent
and therefore contribute no visible product factor.
\item \emph{Multiplier system for half-integer weights.}
Half-integer weights $\tfrac12$ and $\tfrac32$ are realized by
$v_{\eta}^{k}\times v_{H}$, not by the Weil-representation metaplectic
double cover; the integer weight $5$ of $\Delta_{5}$ uses the abelian
multiplier $v_{\eta}^{12}\times v_{H}$ (which is automorphic-trivial
modulo the Maass character $\nu_{\Delta_{5}}$ on the type-II Weyl
generators).  No metaplectic factor enters the Borcherds weight
identity.  See \cite[Theorem~3.1]{GC} for the full classification.
\item \emph{Doubling factor $64^{-1}\Delta_{5}(2Z)$.}
The factor $64^{-1}$ in $\operatorname{den}(\autcor)=64^{-1}
\Delta_{5}(2Z)$ is the same theta-constant $2^{6}$ from
Theorem~\ref{thm:borcherds-product-delta5}, transported through
$Z\mapsto 2Z$; no new constant.  See
Lemma~\ref{lem:appB-doubling}.
\item \emph{The $\kappa_{\mathrm{BKM}}=5$ row.}
The integer $\kappa_{\mathrm{BKM}}=5$ is the row entry of the
Borcherds-weight identity $\operatorname{wt}(F_{j})=c_{N}(0)/2$ in the
catalogue of \cite[Theorem~3.1]{GC}; it agrees with
$\operatorname{wt}(\Delta_{5})=f(0,0)/2=5$ from
\cite[\S2.2]{EZ:1}; it agrees with the entry
$(c_{1}(0),\operatorname{wt})=(10,5)$ in the companion Calabi--Yau
quantum-groups $\kappa$-stratification.
\item \emph{Eichler--Zagier vs Gritsenko--Nikulin convention.}  The
present appendix uses the Eichler--Zagier convention
$\phi_{0,1}=Z_{\mathrm{K3}}/2$ \cite[\S2.2]{EZ:1}; this matches the
Gritsenko--Nikulin Jacobi convention in
\cite[Example~2.4]{GNII} and \cite[\S5.1.1]{GNII} on the
$(n,l)$-coefficient table given in B.1, and matches the singular
theta-lift normalization \cite[Theorem~10.1]{Borcherds95} after the
Mp$_{2}$-vector-valued bridge.  No residual factor of $2$ enters the
exponent $f(nm,l)$ of the Borcherds product.
\end{itemize}

\paragraph{Residual obstructions.}
The four obstruction classes (D0), (O1), (O1)\(^+\), (O2) are not
part of the Borcherds product calculation; they are geometric
content on the source side of $\Delta_{5}$.  Within the automorphic calculation, no
residual Borcherds-product obstruction remains.  The
universal identity $\operatorname{wt}(F_{j})=c_{N}(0)/2$ fixes the row
$j=1$, $\kappa_{\mathrm{BKM}}=5$, and any
disagreement with the Vol~III catalogue, the Borcherds singular
theta-lift weight formula, or the Gritsenko--Cl\'ery classification
would be a contradiction among the stated comparison theorems.

\paragraph{Computations.}
Fourier coefficients are obtained from the rational theta-quotient
expansion; lattice and Gram-pairing identities are the computations of
Appendix~\ref{app:lattice-II21};
the type-II cusp chamber, active support, and Borcherds product are in
Proposition~\ref{prop:bgn-identification-recap}; the doubling identity is
Remark~\ref{rem:bkm-doubling-Z}; the catalogue tables
are in Appendix~\ref{app:eight-finite-constructions}.
